% =========================================================
\chapter{C04 --- Impédances et fonction de transfert}
% =========================================================

\begin{rembox}[title=En bref]
\textbf{De quoi parle ce chapitre ?} Prédire le comportement de circuits R, C (et L) soumis à
des \textbf{sinus}. La notation \textbf{complexe} évite deux difficultés : résoudre des
\textbf{équations différentielles}, et manipuler à la fois des \textbf{amplitudes et des
phases}. \\
\textbf{Les trois objectifs} : la correspondance sinus $\leftrightarrow$ complexes ; le
concept d'\textbf{impédance} $\underline Z$ (C et L se comportent comme une \og résistance
imaginaire\fg{} qui varie avec la fréquence) ; la \textbf{fonction de transfert}, un diviseur
d'impédances qui donne, selon la fréquence, l'atténuation et le déphasage entre entrée et
sortie.
\end{rembox}

\section{Rappels : R, C, L et la difficulté des combinaisons}

\slide{c04-03}{Seuls : facile. Combinés (R et C) : équations différentielles.}

